\documentclass[aspectratio=169,spanish,xcolor=table]{beamer}
\usepackage{iftex}
\ifPDFTeX
  \usepackage[utf8]{inputenc}
  \usepackage[T1]{fontenc}
\fi
\usepackage[spanish,es-tabla,es-noshorthands]{babel}
\usepackage{helvet}
\renewcommand{\familydefault}{\sfdefault}
\usepackage{tikz,pgfplots,array,booktabs,ragged2e}
\pgfplotsset{compat=1.18}
\usetikzlibrary{arrows.meta,positioning,calc,shapes.geometric,fit}

\definecolor{iiiepeNavy}{HTML}{163B5C}
\definecolor{iiiepeBlue}{HTML}{1F6F9F}
\definecolor{iiiepeAqua}{HTML}{20B7A8}
\definecolor{iiiepeOrange}{HTML}{F28C28}
\definecolor{iiiepePaper}{HTML}{F6F8F7}
\definecolor{iiiepeInk}{HTML}{1E2A32}
\definecolor{iiiepePale}{HTML}{EEF5F8}
\definecolor{iiiepeMint}{HTML}{E7F6F4}
\definecolor{iiiepeWarm}{HTML}{FFF3E6}

\usetheme{default}
\usecolortheme{default}
\setbeamertemplate{navigation symbols}{}
\setbeamercolor{normal text}{fg=iiiepeInk,bg=white}
\setbeamercolor{structure}{fg=iiiepeNavy}
\setbeamercolor{frametitle}{fg=white,bg=iiiepeNavy}
\setbeamercolor{block title}{fg=white,bg=iiiepeBlue}
\setbeamercolor{block body}{fg=iiiepeInk,bg=iiiepePale}
\setbeamercolor{item}{fg=iiiepeOrange}

\setbeamertemplate{frametitle}{%
  \nointerlineskip%
  \begin{beamercolorbox}[wd=\paperwidth,ht=0.95cm,dp=0.22cm,leftskip=0.45cm,rightskip=0.45cm]{frametitle}%
    \begin{minipage}[c]{13.2cm}\usebeamerfont{frametitle}\large\bfseries\insertframetitle\end{minipage}%
    \hfill\begin{minipage}[c]{1.6cm}\raggedleft\includegraphics[height=0.56cm]{img/iiiepe/logo-iiiepe-blanco.png}\end{minipage}%
  \end{beamercolorbox}%
}

\setbeamertemplate{footline}{%
  \leavevmode%
  \hbox{%
    \begin{beamercolorbox}[wd=.78\paperwidth,ht=0.5cm,dp=0.18cm,leftskip=0.45cm]{frametitle}%
      \scriptsize\color{white}M. de la Cruz Muñoz · MAEM · Matrícula 00874 · SD26 · Actividad 12 · Sesión 5%
    \end{beamercolorbox}%
    \begin{beamercolorbox}[wd=.22\paperwidth,ht=0.5cm,dp=0.18cm,rightskip=0.45cm]{frametitle}%
      \hfill\scriptsize\color{white}\insertframenumber{} / \inserttotalframenumber%
    \end{beamercolorbox}%
  }%
  \vskip0pt%
}

\newcommand{\Pill}[2]{%
  \tikz[baseline=(x.base)]{\node[rounded corners=3pt,fill=#1!14,draw=#1!75!black,thick,inner xsep=6pt,inner ysep=4pt,align=center] (x) {\scriptsize\bfseries #2};}%
}

\title{Escuela Tradicional vs. NEM y Tres Reportes de Lectura}
\subtitle{Actividad 12 · Sesión 5 (Meirieu 2020 · Sawyer 2008 · Martínez Rizo 2026)}
\author{Mtro. Martín Jonathan de la Cruz Muñoz · Matrícula 00874}
\institute{IIIEPE · Maestría en Aprendizaje y Enseñanza de las Matemáticas · Grupo A}
\date{19 de noviembre de 2026}

\begin{document}

\begin{frame}[plain]
  \begin{tikzpicture}[remember picture,overlay]
    \fill[iiiepeNavy] (current page.north west) rectangle ([xshift=10.3cm]current page.south west);
    \fill[iiiepeBlue] ([xshift=10.3cm]current page.north west) rectangle (current page.south east);
    \fill[iiiepeOrange] ([xshift=0.8cm,yshift=-6.45cm]current page.north west) rectangle ([xshift=9.3cm,yshift=-6.53cm]current page.north west);
    \node[anchor=north west,text width=8.7cm,align=left] at ([xshift=0.8cm,yshift=-1.15cm]current page.north west) {%
      {\color{iiiepeAqua!45!white}\small\bfseries MAESTRÍA EN APRENDIZAJE Y ENSEÑANZA DE LAS MATEMÁTICAS · SD26\par}%
      \vspace{0.35cm}
      {\color{white}\Large\bfseries \inserttitle\par}%
      \vspace{0.25cm}
      {\color{white!90}\normalsize \insertsubtitle\par}%
    };
    \node[anchor=south west,text width=8.7cm,align=left] at ([xshift=0.8cm,yshift=0.75cm]current page.south west) {%
      {\color{white}\normalsize\bfseries Mtro. Martín Jonathan de la Cruz Muñoz\par}%
      {\color{white!85}\scriptsize Matrícula 00874 · martin.delacruz@campus.iiiepe.edu.mx\par}%
      {\color{iiiepeAqua!35!white}\scriptsize Monterrey, Nuevo León · 19 de noviembre de 2026\par}%
    };
    \node[anchor=center,text width=5.0cm,align=center] at ([xshift=13.15cm,yshift=-4.5cm]current page.north west) {%
      \includegraphics[width=3.6cm]{img/iiiepe/logo-iiiepe-blanco.png}\par\vspace{0.45cm}%
      {\color{white}\large\bfseries IIIEPE\par}%
      \vspace{0.1cm}%
      {\color{white!85}\scriptsize Fundamentos para la enseñanza\\y el aprendizaje I\par}%
    };
  \end{tikzpicture}
\end{frame}
\begin{frame}{Contraste Estructural: Escuela Tradicional vs. NEM}
  \renewcommand{\arraystretch}{1.25}
  \scriptsize
  \begin{tabular}{p{2.6cm}p{5.2cm}p{5.5cm}}
    \toprule
    \textbf{Dimensión} & \textbf{Escuela tradicional (Scallon, 2004)} & \textbf{Nueva Escuela Mexicana (SEP, 2022)} \\
    \midrule
    Finalidad & Reproducción memorística de hechos. & Formación humanista, crítica y comunitaria. \\
    Currículo & Asignaturas aisladas fragmentadas. & Campos formativos y programa analítico. \\
    Problemas & Ejercicios cerrados artificiales. & Situaciones auténticas con múltiples fuentes. \\
    Evaluación & Prueba puntual unidimensional. & Proceso formativo, retroalimentación y diálogo. \\
    Error y equidad & Sanción individual del rezago. & Diagnóstico y apoyos para derecho común. \\
    \bottomrule
  \end{tabular}
\end{frame}

\begin{frame}{Tres Escalas Complementarias de Mejora Educativa}
  \begin{columns}[T]
    \column{0.31\textwidth}
      \begin{block}{1. Philippe Meirieu (2020)}
        \scriptsize\textbf{Escala institucional:}\\
        Construir aula común (\textit{yo + nosotros}), superar taylorismo y ofrecer «algo más» frente al conductismo digital.
      \end{block}
    \column{0.31\textwidth}
      \begin{block}{2. Keith Sawyer (2008)}
        \scriptsize\textbf{Escala de aula y cognición:}\\
        Aprendizaje profundo, saberes previos, articulación, metacognición y andamiaje progresivo.
      \end{block}
    \column{0.31\textwidth}
      \begin{block}{3. Martínez Rizo (2026)}
        \scriptsize\textbf{Escala sistémica (PISA 2025):}\\
        Lectura histórica ($409\to 395\to 388$ en matemáticas), cobertura y sistema nacional propio de evaluación.
      \end{block}
  \end{columns}
\end{frame}

\begin{frame}{De la Ecuación Rutinaria al Análisis Formativo en Matemáticas}
  \begin{block}{Ejemplo comparado: $0.5x + 2 = 0.2x + 3.8 \quad\text{y}\quad C_A(x)=30+8x \text{ vs. } C_B(x)=12x$}
    \begin{itemize}\normalsize
      \item \textbf{Tradicional:} Despejar $x=6$ o $x=7.5$ sin significado físico ni verificación.
      \item \textbf{Comprensivo (NEM + Sawyer + Meirieu):} Justificar equivalencia algebraica, vincular tabla--gráfica--ecuación, dar andamiaje con ejes guía y devolver errores como hipótesis revisables.
    \end{itemize}
  \end{block}
\end{frame}

\begin{frame}{Conclusiones y Referencias (APA 7)}
  \begin{itemize}\small
    \item La mejora educativa requiere alinear propósito público (Meirieu), diseño cognitivo (Sawyer) y evidencia contextualizada (Martínez Rizo).
    \item La NEM aporta un marco integrador cuando evita reducir el proyecto o la evaluación a mero cumplimiento formal.
  \end{itemize}
  \vspace{0.3cm}
  \scriptsize
  \textbf{Referencias:} Martínez Rizo, F. (2026). Resultados de PISA 2025. \textit{Nexos}; Meirieu, P. (2020). \textit{MCEP}; Sawyer, R. K. (2008). \textit{OECD}; Scallon, G. (2004); SEP (2022). \textit{Plan de Estudio}.
\end{frame}

\end{document}
